\documentclass[11pt,letterpaper]{article}

\usepackage[letterpaper,margin=0.55in,top=0.40in,bottom=0.38in]{geometry}
\usepackage[T1]{fontenc}
\usepackage{newtxtext}
\usepackage{helvet}
\usepackage{hyperref}
\usepackage{enumitem}
\usepackage{xcolor}
\usepackage{array}
\usepackage{tabularx}
\usepackage{microtype}
\usepackage{ragged2e}

\definecolor{linkblue}{RGB}{0,102,204}
\hypersetup{
  colorlinks=true,
  urlcolor=linkblue
}

\pagestyle{empty}
\setlength{\parindent}{0pt}
\setlength{\tabcolsep}{0pt}
\setlength{\parskip}{0pt}
\setlength{\emergencystretch}{3em}
\hyphenpenalty=10000
\exhyphenpenalty=10000
\linespread{0.99}

\newcommand{\namestyle}{\fontfamily{phv}\selectfont}
\newcommand{\sectionstyle}{\fontfamily{phv}\bfseries\selectfont}
\newcommand{\link}[2]{\href{#1}{\textcolor{linkblue}{\underline{#2}}}}
\newcolumntype{Y}{>{\raggedright\arraybackslash}X}

\newcommand{\resumesection}[1]{%
  \vspace{2pt}%
  {\centering{\sectionstyle\large #1}\par}%
  \nointerlineskip\vspace{2pt}\hrule\vspace{3pt}%
}

\newcommand{\educationentry}[3]{%
  \begin{tabularx}{\linewidth}{@{}Yr@{}}
    {\sectionstyle #1} \textbar{} \textit{#2} & #3
  \end{tabularx}%
}

\newcommand{\experienceentry}[5]{%
  \begin{tabularx}{\linewidth}{@{}Yr@{}}
    {\sectionstyle #1} \textbar{} \textit{#2} \textbar{} #3 & #4 \textbar{} #5
  \end{tabularx}%
}

\newcommand{\projectentry}[4]{%
  \begin{tabularx}{\linewidth}{@{}Yr@{}}
    \link{#1}{\sectionstyle #2} \textbar{} #3 & #4
  \end{tabularx}%
}

\newcommand{\projectentrynolink}[3]{%
  \begin{tabularx}{\linewidth}{@{}Yr@{}}
    {\sectionstyle #1} \textbar{} #2 & #3
  \end{tabularx}%
}

\newenvironment{resumebullets}{%
  \RaggedRight
  \begin{itemize}[
    label=\raisebox{0.25ex}{\rule{3pt}{3pt}},
    leftmargin=0.32in,
    labelsep=0.12in,
    itemsep=0pt,
    topsep=0pt,
    parsep=0pt,
    partopsep=0pt
  ]%
}{\end{itemize}}

\begin{document}

\begin{minipage}[t]{0.28\linewidth}
  \raggedright
  \link{http://snehithn.com/}{http://snehithn.com/}\\
  \link{mailto:snehithn5@gmail.com}{snehithn5@gmail.com}
\end{minipage}\hfill%
\begin{minipage}[t]{0.38\linewidth}
  \centering
  \vspace{-5pt}
  {\namestyle\fontsize{25}{29}\selectfont Snehith Nayak}
\end{minipage}\hfill%
\begin{minipage}[t]{0.28\linewidth}
  \raggedleft
  SF Bay Area, CA
\end{minipage}

\resumesection{EDUCATION}

\educationentry{University of Michigan Ann Arbor}{M.S. Applied Data Science}{Aug 2025 - Present}
\educationentry{University of California Santa Barbara}{B.S. Computer Engineering}{Jun 2024}

\resumesection{WORK EXPERIENCE}

\experienceentry{Applied Materials}{Software Engineer \link{https://www.appliedmaterials.com/us/en/semiconductor/solutions-and-software/ai-x.html}{AIx}}{Distributed Systems, ML Infra}{Santa Clara, CA}{July 2024 - Present}
\begin{resumebullets}
  \item Engineered a real-time data platform on Kubernetes, Apache Kafka (AWS MSK), and ClickHouse/DuckDB,
  ingesting \textasciitilde35M rows/month of fab metrology and process data into a multi-billion-row store powering
  search, analytics, and recommendations across process-development experiments.
  \item Took the platform 0$\to$1 to production (Feb 2025) and scaled to 600+ monthly active users, hardening
  the pipeline with priority-tiered Kafka topics, dead-letter queues, and circuit breakers for fault isolation.
  \item Built a recommendation layer that analyzes linked experiment and metrology data to suggest which process
  parameters engineers should tune, providing data-driven starting points for new semiconductor process recipes.
  \item Designed an experiment-tracking and metadata-logging system that makes every fab process and metrology
  run reproducible and directly comparable by configuration and result.
\end{resumebullets}

\experienceentry{KLA}{Machine Learning Engineer Intern - \link{https://snehithn.com/images/KLA-BBP.pdf}{BBP}}{Python, Image Classification}{Milpitas, CA}{June-Sept 2023}
\begin{resumebullets}
  \item Developed a machine-learning autoencoder using TensorFlow for BBP Nuisance Filtering, reducing
  cluster count by 30\% in IBM CAD and Samsung GMK dataset compared to DBG2 (current implementation).
  \item Engineered an advanced, adaptable framework using convolutional layers, optimizing defect detection
  and classification across multiple datasets, specialized by training with unlabeled BBP hotspot images.
  \item Utilized HDBSCAN for efficient clustering of latent features, validated by silhouette scores.
\end{resumebullets}

\experienceentry{Synopsys}{Application Development Intern}{Java, SQL, Data Processing}{Santa Clara, CA}{June-Sept 2022}
\begin{resumebullets}
  \item Automated contractor onboarding process saving 2,200 employee hours annually, improving system
  efficiency for over 2,000 users.
\end{resumebullets}

\resumesection{PROJECTS}

\projectentry{https://pixelpane.app}{Pixel Pane -- AI Assistant}{SwiftUI, MLX, Cloudflare Workers, Anthropic API}{May 2026}
\begin{resumebullets}
  \item Built model-agnostic SwiftUI macOS AI assistant supporting any MLX-compatible backend, Anthropic Claude, and Apple Foundation Models — with per-provider health detection and lifecycle management.
  \item Deployed Cloudflare Workers proxy normalizing Anthropic streaming into SSE with per-user rate limiting, KV quota enforcement, and cryptographic request tracing.
  \item Engineered extensible action dispatch system with heuristic content classification, built-in handlers for translation and code debugging, and user-defined side effects.
  \item Designed permission-gated security model with explicit file grants and scoped side-effect authorization before any system-level operation.
\end{resumebullets}

\projectentry{https://isocode.dev}{iso-code -- Safe Worktree Management for AI Agents}{Rust, CLI, MCP Server, Git}{Apr 2026}
\begin{resumebullets}
  \item Built a Rust library, CLI, and MCP server providing safe git worktree lifecycle management for AI coding
  agents across Claude Code, Cursor, and Copilot.
  \item Prevented data loss from unmerged-commit deletion, orphaned worktrees, and git-crypt corruption with
  atomic state writes and cleanup-on-failure guarantees.
\end{resumebullets}

\projectentry{https://snehithn.com/images/CAPSTONE.pdf}{Health Monitoring Wearable -- Capstone}{Python, C, Java, ASIC Design}{Nov 2023-June 2024}
\begin{resumebullets}
  \item Designed and manufactured a specialized health monitoring wearable for nursing home residents,
  focusing on comprehensive remote monitoring and enhanced data accuracy.
  \item Incorporated advanced sensors in the wearable for vital signs, activity tracking, temperature monitoring,
  ambient noise detection, and continuous heart rate and oxygen level measurement.
  \item Engineered a data-efficient Android app for real-time health data visualization and trend analysis.
  \item Designed and manufactured a 4-Layer PCB design in Altium.
\end{resumebullets}

\resumesection{SKILLS}

{\sectionstyle Languages:} Python, SQL, C++\\
{\sectionstyle ML \& Data:} PyTorch, TensorFlow, ClickHouse, DuckDB, PyArrow\\
{\sectionstyle Infrastructure:} Kubernetes, Kafka, AWS (MSK, S3, CloudWatch), Redis, Docker, FastAPI, Git\\
{\sectionstyle AI tooling:} MCP, Claude API, OpenAI API\\

\end{document}
